% 4.3 章 · 守恒律与有限体积
% 这一章的主线：积分形式 -> 弱解 -> Rankine–Hugoniot -> 熵条件 -> Godunov 格式
%   -> 有限体积的望远镜求和 -> Euler 方程组的 Riemann 解扇。
% 每个纸面结论都配本章断言的实测数字（tests/4.3/ 的断言报告）。
% 断言附表由 tools/build_book.py 生成，\input 进来，不手抄。

\numlabchapter{守恒律与有限体积}{4.3}

\section{为什么守恒律必须从积分形式开始}

标量守恒律
\begin{equation}
  u_t+f(u)_x=0,\qquad x\in(0,L),\ t>0,\qquad u(x,0)=u_0(x)
  \label{eq:cl}
\end{equation}
看起来只是平流方程把 $a u_x$ 换成了 $f(u)_x$，但两者的数学性质有一个决定性的差别：
\eqref{eq:cl} 的\textbf{经典解会在有限时间破裂}。Burgers 方程 $f(u)=u^2/2$ 是最小的例子：
特征线 $x(t)=x_0+t\,u_0(x_0)$ 是直线，斜率就是初值本身；初值递减的地方，
后面的特征线追上前面的，解出现多值。

物理量（质量、动量、能量）在破裂之后仍然有定义，所以正确的出发点是积分形式。
对任意区间 $[a,b]$，
\begin{equation}
  \frac{d}{dt}\int_a^b u\,dx
  = f(u(a,t))-f(u(b,t)).
  \label{eq:integral}
\end{equation}
\eqref{eq:integral} 只要求 $u$ 是 $L^1_{\text{loc}}$，不要求它可导；它把 ``$u_t+f(u)_x=0$''
从一个逐点等式降格成一个\textbf{分布意义}下的等式。这一降格不是为了偷懒：
破裂后的熵解本来就不是经典解，任何只对光滑解成立的推导都管不到它。

\begin{definition}[弱解]
若 $u\in L^\infty(\mathbb{R}\times[0,\infty))$ 且对一切 $\varphi\in C_c^\infty(\mathbb{R}\times[0,\infty))$
\begin{equation}
  \int_0^\infty\!\!\int_{\mathbb{R}}
  \bigl(u\,\varphi_t+f(u)\,\varphi_x\bigr)\,dx\,dt
  +\int_{\mathbb{R}}u_0(x)\,\varphi(x,0)\,dx=0,
  \label{eq:weak}
\end{equation}
则称 $u$ 是~\eqref{eq:cl} 的弱解。
\end{definition}

弱解不一定唯一。取 $u_L=0$、$u_R=1$ 的 Burgers 数据：一条右行激波（速度 $1/2$）
与一条左行膨胀激波（速度 $1/2$，方向相反）都满足~\eqref{eq:weak}。
所以``弱解''这个词本身只是把可选项扩大了；要选出物理上那一个，需要熵条件。
这是本章所有``防误导''讨论的源头：\textbf{满足守恒律的数值解可以收敛到错误的东西}。

\section{有限体积：守恒性是一条恒等式，不是一句口号}

\begin{definition}[单元平均与数值通量]
把区间分成 $N$ 个长度 $\Delta x$ 的单元 $I_j=[x_{j-1/2},x_{j+1/2}]$，
令 $U_j^n\approx\frac1{\Delta x}\int_{I_j}u(x,t_n)\,dx$。有限体积格式写成
\begin{equation}
  U_j^{n+1}=U_j^n-\frac{\Delta t}{\Delta x}
  \bigl(F_{j+1/2}-F_{j-1/2}\bigr),
  \qquad F_{j+1/2}=F(U_j,U_{j+1}).
  \label{eq:fv}
\end{equation}
$F$ 是\textbf{数值通量}：一致（$F(u,u)=f(u)$）、单调或至少满足某种熵估计。
\end{definition}

\begin{theorem}[离散守恒]
\label{thm:conservation}
设~\eqref{eq:fv} 在 $j=1,\dots,N$ 上成立，$F_{1/2}$ 与 $F_{N+1/2}$ 是边界通量，则
\begin{equation}
  \sum_{j=1}^N U_j^{n+1}\Delta x
  =\sum_{j=1}^N U_j^{n}\Delta x
  -\Delta t\bigl(F_{N+1/2}-F_{1/2}\bigr).
  \label{eq:conservation}
\end{equation}
特别地，周期边界（$F_{1/2}=F_{N+1/2}$）下总质量逐位不变。
\end{theorem}

\begin{proof}
把~\eqref{eq:fv} 乘以 $\Delta x$ 并对 $j$ 求和。右端的内部通量成对出现：
$F_{j+1/2}$ 在单元 $j$ 的方程里取负号，在单元 $j+1$ 的方程里取正号，
求和后只剩两个端点。边界条件只通过 $F_{1/2}$ 与 $F_{N+1/2}$ 进入。
\end{proof}

这条证明短得不像一个定理，但它正是本章最重要的结构性事实：
\textbf{内部通量怎么错都不破坏守恒}。迎风格式、Rusanov、Osher、MUSCL，
甚至通量算错的代码，只要写成~\eqref{eq:fv} 并且边界通量记账一致，总质量就严格守恒。
反过来，\textbf{只有不写成通量差的离散才会漏}——比如把 $u_t+f(u)_x=0$ 改写成
特征形式 $u_t+u\,u_x=0$ 再对 $u\,u_x$ 做中心差分或迎风差分。
两种写法在光滑区逐点相同，在间断处完全不同。

本章的断言把这件事量成了两个数。取 $u_L=1$、$u_R=0$ 的激波、$N=160$、$t=0.6$、
出流边界，三套格式用\textbf{同一套记账规则}：
\begin{center}
\begin{tabular}{@{}lrr@{}}
\toprule
格式 & 最大账本残差 $|M-M_0+(\text{out}-\text{in})|$ & 末态质量 \\
\midrule
一阶 Godunov（守恒） & $5.83\times10^{-16}$ & $0.800000$ \\
MUSCL--Hancock（守恒） & $1.17\times10^{-15}$ & $0.800000$ \\
特征形式（非守恒） & $3.00\times10^{-1}$ & $0.500000$ \\
\bottomrule
\end{tabular}
\end{center}
出流边界下质量本来就该从 $0.5$ 变到 $0.8$（累计流入 $0.3$），所以
\textbf{质量曲线本身不是守恒判据}；残差才是。三者的差有 15 个数量级，
而解在光滑区看起来完全一样——这是本章第一个、也是最容易被忽略的陷阱。

\section{Rankine–Hugoniot：激波速度从守恒律里唯一确定}

\begin{theorem}[Rankine--Hugoniot 条件]
\label{thm:rh}
设 $u$ 是分片光滑弱解，在曲线 $x=s t$ 两侧的极限为 $u_L,u_R$。则
\begin{equation}
  s\,(u_R-u_L)=f(u_R)-f(u_L),
  \qquad\text{即}\qquad
  s=\frac{f(u_R)-f(u_L)}{u_R-u_L}.
  \label{eq:rh}
\end{equation}
\end{theorem}

\begin{proof}
取一个跨过间断的矩形控制体 $[-\varepsilon+s t,\ \varepsilon+s t]\times[t_1,t_2]$，
在其上对~\eqref{eq:integral} 积分。让 $\varepsilon\to0^+$：
体积项给出 $\int_{t_1}^{t_2}s(u_R-u_L)\,dt$（上下边界的贡献因高度趋零而消失），
通量项给出 $\int_{t_1}^{t_2}(f(u_R)-f(u_L))\,dt$。对任意区间成立，
故被积函数相等。
\end{proof}

对 Burgers，$s=(u_L+u_R)/2$。取 $u_L=1,u_R=0$，$s=1/2$。
本章用 $n=400$ 的一阶 Godunov 真推进到 $t=0.5$，追踪 $u=1/2$ 的等值线位置，
再做位置-时间线性拟合：$\hat s=0.50003848$（$r^2=0.999999$），相对误差 $7.7\times10^{-5}$。
误差主要来自激波被抹宽后等值线位置的离散抖动，而不是格式的相速度。

定理~\ref{thm:rh} 只用了守恒，没有用熵。这就是为什么``RH 满足''不能作为
正确性的判据：膨胀激波也满足它。

\section{熵条件：在 RH 之外再选一次}

\begin{definition}[Lax 熵条件；凸通量]
对凸通量（$f''>0$），跨激波
\begin{equation}
  f'(u_L)>s>f'(u_R).
  \label{eq:lax}
\end{equation}
对 Burgers 就是 $u_L>s>u_R$，即\textbf{左状态跑得比激波快、右状态跑得慢}。
若 $u_L<u_R$，~\eqref{eq:rh} 给出的曲线两侧特征线是发散的，
它是一条\textbf{膨胀激波}：数学上满足弱解定义，物理上不存在。
\end{definition}

\begin{theorem}[熵不等式]
\label{thm:entropy}
设存在凸熵 $\eta(u)$ 与熵通量 $q(u)$，$q'=\eta'f'$。物理弱解必须满足分布意义下的
\begin{equation}
  \eta(u)_t+q(u)_x\le 0.
  \label{eq:entropy}
\end{equation}
对 Burgers，取 $\eta=u^2/2$、$q=u^3/3$：光滑区~\eqref{eq:entropy} 取等号，
激波处严格小于零（熵产为正），膨胀激波则给出严格大于零。
\end{theorem}

\begin{proof}[证明思路]
在方程里加一个消失的黏性项 $\varepsilon u_{xx}$，对光滑解乘以 $\eta'(u)$ 得
$\eta_t+q_x=\varepsilon\,\eta''(u)\,u_x^2\le0$（$\eta$ 凸）。令 $\varepsilon\to0^+$，
黏性解收敛到满足~\eqref{eq:entropy} 的弱解；而膨胀激波作为消失黏性的极限不出现。
\end{proof}

本章的数值证据来自跨声速稀疏波 $u_L=-1$、$u_R=1$：
\begin{itemize}
  \item 精确 Godunov 通量取临界点 $u_s=0$，通量 $f(u_s)=0$；
  \item 朴素迎风通量按最大波速取一侧，给 $f(u_R)=1/2$。
\end{itemize}
两者相差 $0.5$，\textbf{不是舍入，是两种不同的弱解}。
真推进到 $t=0.3$：Godunov 的末态熵产 $-5.36\times10^{-3}$（正确，熵被耗散），
朴素迎风 $+2.00\times10^{-1}$（错误，熵被凭空制造）。
与此同时，两者的总变差都没有增加、RH 都满足、格式都单调——
\textbf{常规检查全部通过}。这就是本章第二个陷阱：熵条件是唯一能抓住它的判据。

\section{Godunov 格式与 Riemann 问题}

\begin{definition}[Godunov 通量]
在每个界面 $x_{j+1/2}$ 上解精确 Riemann 问题（左状态 $U_j$、右状态 $U_{j+1}$），
取相似变量 $\xi=x/t=0$ 处的数值通量：
\begin{equation}
  F^{\text{Gud}}_{j+1/2}=f\bigl(U^\star(0;U_j,U_{j+1})\bigr).
\end{equation}
\end{definition}

对 Burgers，$\xi=0$ 处的 Riemann 解有三种情形：
\begin{itemize}
  \item $u_L>u_R$：激波，$u^\star(0)=u_L$（因为 $s>0$ 时左状态占 $x=0$ 的上风侧），$F=f(u_L)$；
  \item $u_L\le u_R$ 且 $0\notin[u_L,u_R]$：单侧状态，$F=f(u_L)$ 或 $f(u_R)$；
  \item 跨声速（$u_L<0<u_R$）：稀疏波扇面覆盖 $x=0$，$u^\star(0)=0$，$F=f(0)=0$。
\end{itemize}
最后一条是 Godunov 格式与朴素迎风通量的全部差别：迎风只看``风往哪边吹''，
Godunov 还看``$x=0$ 处到底是哪个状态''。

\begin{theorem}[Godunov 格式的熵稳定性]
\label{thm:godunov-entropy}
若 $F$ 是精确 Riemann 通量，则格式满足离散熵不等式
\begin{equation}
  \eta(U_j^{n+1})\le\eta(U_j^n)
  -\frac{\Delta t}{\Delta x}\bigl(Q_{j+1/2}-Q_{j-1/2}\bigr),
\end{equation}
其中 $Q$ 由 Riemann 扇面的熵通量给出。特别地，收敛的 Godunov 解满足~\eqref{eq:entropy}。
\end{theorem}

\begin{proof}[证明思路]
精确 Riemann 解的熵产在每条激波上非负、在稀疏波内为零（等熵），
所以每个界面贡献的数值熵通量满足所需的单侧估计。沿 $t$ 积分并用
单元平均的凸性（Jensen 不等式）即可得到离散不等式。
\end{proof}

本章用 $n=200$、$\text{cfl}=0.4$ 验证了 Godunov 的 TVD：
总变差从 $1.0000$ 到 $1.0000$，逐点最大单步增量为
$3.3\times10^{-16}$（舍入级）。注意 TVD 只说明``不产生新振荡''，
它对熵条件保持沉默。

\section{Euler 方程组的 Riemann 解扇}

Euler 方程组
\begin{equation}
  \mathbf U=\begin{pmatrix}\rho\\ \rho u\\ E\end{pmatrix},
  \qquad
  \mathbf F=\begin{pmatrix}\rho u\\ \rho u^2+p\\ u(E+p)\end{pmatrix},
  \qquad
  E=\frac{p}{\gamma-1}+\frac12\rho u^2
  \label{eq:euler}
\end{equation}
的 Riemann 解比标量情形多了一种波：接触间断。它的结构是
\begin{equation}
  \underbrace{\text{左波}}_{\text{激波或稀疏波}}\quad
  \underbrace{\text{接触间断}}_{u,p\ \text{连续},\ \rho\ \text{跳}}\quad
  \underbrace{\text{右波}}_{\text{激波或稀疏波}} .
  \label{eq:fan}
\end{equation}

\begin{theorem}[星区压强方程]
\label{thm:pstar}
令 $f_L(p)$、$f_R(p)$ 为左/右波的 ``压强函数''（激波用 RH、稀疏波用等熵关系），
则星区压强 $p^\star$ 是方程
\begin{equation}
  f_L(p^\star)+f_R(p^\star)+u_R-u_L=0
  \label{eq:pstar}
\end{equation}
的唯一正根；星区速度 $u^\star$ 由任一$f$给出。
\end{theorem}

\begin{proof}[证明思路]
$f_L,f_R$ 在 $(0,\infty)$ 上连续、单调递增（激波段 $f'>0$，稀疏波段由
$f_K(p)=\frac{2c_K}{\gamma-1}\bigl[(p/p_K)^{(\gamma-1)/(2\gamma)}-1\bigr]$ 也递增），
且在 $p\to0^+$ 与 $p\to\infty$ 时分别趋于有限值与 $+\infty$，故和函数严格递增；
中间值定理给出唯一根。若 $f_L(0)+f_R(0)+u_R-u_L>0$，则出现真空，
本章的三个算例都避开这一支。
\end{proof}

用二分解出 Sod 的星区值，与 Toro 的教科书值对拍，最大相对差
$6.4\times10^{-9}$：
\[
  p^\star=0.30313018,\qquad u^\star=0.92745262,\qquad
  \rho^\star_L=0.42631943,\qquad \rho^\star_R=0.26557371 .
\]
跨右激波的三条 RH 条件（质量、动量、能量，按跳跃尺度归一）残差为
$9.2\times10^{-16}$、$2.3\times10^{-16}$、$1.2\times10^{-15}$。
接触间断两侧在 $\xi=u^\star$ 的 $10^{-8}$ 邻域内 $|\Delta p|=|\Delta u|=0$，
而 $|\Delta\rho|=0.161$——这正是~\eqref{eq:fan} 的结构。

\begin{theorem}[HLLC 精确分辨静止接触]
\label{thm:hllc-contact}
若 $u_L=u_R=0$、$p_L=p_R=p$、$\rho_L\ne\rho_R$，则 HLLC 通量为
\begin{equation}
  \mathbf F^{\text{HLLC}}=(0,\ p,\ 0)^{\mathsf T},
\end{equation}
与密度无关。
\end{theorem}

\begin{proof}
静止接触的中间波速 $S_M=0$；HLLC 的星区选择退化为在接触两侧各取常状态，
而两侧的动量为零、总能量只通过 $p$ 进入通量，故动量通量就是 $p$，
质量与能量通量为零。
\end{proof}

本章断言的实测值是 $[0,\ 1.000000000000,\ 0]$（误差 $\le10^{-14}$）。
这条性质说明：\textbf{接触间断是否被抹宽，是检验数值通量分辨率的直接读数}；
一阶格式把它抹成 $O(\sqrt{\Delta x})$ 宽，二阶格式 $O(\Delta x)$。

\section{特征线与破散时间}

对 Burgers，$f'(u)=u$，特征线是直线
\begin{equation}
  x(t)=x_0+t\,u_0(x_0).
  \label{eq:char}
\end{equation}
相邻两条特征线（$u_i>u_{i+1}$）的交点
\begin{equation}
  t_\ast=\frac{x_{i+1}-x_i}{u_i-u_{i+1}},
  \qquad
  x_\ast=x_i+t_\ast u_i,
  \label{eq:cross}
\end{equation}
最早的交点就是破散时间 $t_b=-1/\min_x u_0'(x)$。
对 $u_0=A\sin(2\pi x)$，$\min u_0'=-2\pi A$，故
\begin{equation}
  t_b=\frac{1}{2\pi A}.
  \label{eq:tb}
\end{equation}

本章用三条独立路线量同一个 $t_b$：
闭式 $0.3183098862$、20001 点扫描 $0.3183098862$（逐位相同，因为极值点落在网格上）、
40 条特征线的首个交点 $0.3186373711$（相对差 $1.03\times10^{-3}$，
误差来自离散起点 $x_0=(i+1/2)/n$ 与真实极值点的偏移）。
三条路线一致到 $O(1/n)$，这就是``理论值''在数值实验里的正确用法：
\textbf{允许它比网格精度更准，但不允许它和网格无关}。

破散之后（$t>t_b$），\eqref{eq:char} 的回溯不再唯一：同一条 $x$ 对应三条特征线，
``最靠左的一支''只是多值解的一个分支。拿它当精确解算误差，误差会随 $\Delta x$
漂亮地收敛——但收敛到一个\textbf{不是熵解}的东西。这是本章第三个陷阱，
断言里用 \texttt{exact\_valid} 掩码把这种比较整条作废。

\section{为什么二阶不会自动更守恒}

MUSCL--Hancock 在单元内做斜率重构，再用半步预测更新界面状态：
\begin{equation}
  U_j^{L}=U_j^n+\frac12\,s_j,\qquad
  U_j^{R}=U_j^n-\frac12\,s_j,\qquad
  s_j=\operatorname{minmod}\bigl(U_j-U_{j-1},\ U_{j+1}-U_j\bigr).
  \label{eq:muscl}
\end{equation}
$s_j=0$ 时它逐位退化成 Godunov（限制器把极值处降成一阶）。
限制器保证 TVD，代价是在极值处损失二阶精度；这正是 Godunov 定理
（线性单调格式至多一阶）允许的折中。

本章实测：一阶与 MUSCL 的账本残差分别是 $5.83\times10^{-16}$ 与
$1.17\times10^{-15}$——两者都在舍入级，MUSCL 只是大了 2 倍。
\textbf{``更准''与``更守恒''是两条独立的账}：二阶重构没有改变
通量差的望远镜结构，所以守恒性不随精度阶变化；它改变的是通量的数值，
从而改变解的形状、熵产和激波位置。

\section{本章的实测数字一览}

\begin{center}
\begin{tabular}{@{}lll@{}}
\toprule
量 & 实测 & 理论/判据 \\
\midrule
Burgers 激波速度（n=400） & $0.50003848$ & RH $=0.5$，相对差 $7.7\times10^{-5}$ \\
Sod $p^\star$ & $0.30313018$ & Toro $0.30313018$，相对差 $6.4\times10^{-9}$ \\
Sod $u^\star$ & $0.92745262$ & Toro $0.92745262$ \\
跨右激波 RH 最大残差 & $1.19\times10^{-15}$ & 舍入级 \\
接触两侧 $|\Delta p|/|\Delta u|/|\Delta\rho|$ & $0/0/0.161$ & 连续/连续/跳跃 \\
静止接触的 HLLC 通量 & $(0,1,0)$ & 精确 \\
跨声速稀疏波 Godunov / 迎风通量 & $0$ / $0.5$ & 临界点 / 单侧 \\
Godunov 熵产（稀疏波，t=0.3） & $-5.36\times10^{-3}$ & $\le0$ \\
迎风熵产（同一算例） & $+2.00\times10^{-1}$ & $>0$，熵违背 \\
守恒账本残差（Godunov/MUSCL/迎风） & $5.83\times10^{-16}/1.17\times10^{-15}/3.0\times10^{-1}$ & 守恒/守恒/漏水 \\
破散时间：闭式/扫描/特征线 & $0.3183098862/0.3183098862/0.3186373711$ & 相对差 $\le1.03\times10^{-3}$ \\
\bottomrule
\end{tabular}
\end{center}

% 本文件由 tools/build_book.py 从 dist/reports/assertions.json 生成，请勿手改。
\begin{longtable}{@{}p{0.06\textwidth}p{0.32\textwidth}p{0.22\textwidth}p{0.16\textwidth}@{}}
\caption{章节 4.3 数值断言（生成时间 2026-10-19T07:20:49-0400）}\\
\toprule
类别 & 断言 & 理论值 & 实测值 \\
\midrule
\endfirsthead
\toprule
类别 & 断言 & 理论值 & 实测值 \\
\midrule
\endhead
\bottomrule
\endlastfoot

A & Godunov 实测激波速度(n=400) & RH: s=\allowbreak{}0.5 & 0.500038 \checkmark \\
A & 同一 xi 在两个时刻的精确解最大差 & 自相似 -> 0 & 0 \checkmark \\
A & Sod 星区值 vs Toro 教科书 & p*=\allowbreak{}0.30313018 等 & $6.431\times 10^{-9}$ \checkmark \\
A & Sod 右激波三条 RH 残差的最大值 & <=\allowbreak{} 1e-13(\allowbreak{}双向归一) & $1.187\times 10^{-15}$ \checkmark \\
A & 接触间断两侧：|$\Delta$p|, |$\Delta$u|, |$\Delta$rho| & dp=\allowbreak{}du=\allowbreak{}0，drho>0.1 & dp=\allowbreak{}0.0e+00,\allowbreak{} du=\allowbreak{}0.0e+00,\allowbreak{} drho=\allowbreak{}0.161 \checkmark \\
A & 静止接触的 HLLC 质量/动量/能量通量 & [0,\allowbreak{} p,\allowbreak{} 0] & [0.0e+00,\allowbreak{} 1.000000000000,\allowbreak{} 0.0e+00] \checkmark \\
A & 跨声速稀疏波的 Godunov vs upwind 通量 & Godunov=\allowbreak{}0,\allowbreak{} upwind=\allowbreak{}0.5 & godunov=\allowbreak{}0.0,\allowbreak{} upwind=\allowbreak{}0.5 \checkmark \\
A & 破散时间闭式/扫描/特征线交点 & 闭式 0.318309886184 & scan=\allowbreak{}0.318309886,\allowbreak{} lines=\allowbreak{}0.318637371 \checkmark \\
B & 三本账的最大残差(godunov / muscl / advective) & 守恒两条 \textasciitilde{}1e-16，非守恒 \textasciitilde{}0.3 & 5.83e-16 /\allowbreak{} 1.17e-15 /\allowbreak{} 3.00e-01 \checkmark \\
B & Godunov 的总变差最大单步增量 & <=\allowbreak{} 1e-12(\allowbreak{}TVD) & $3.331\times 10^{-16}$ \checkmark \\
B & 跨声速稀疏波末态熵产(Godunov / upwind) & Godunov <=\allowbreak{} 0，upwind > 0 & godunov=\allowbreak{}-5.358e-03,\allowbreak{} upwind=\allowbreak{}2.000e-01 \checkmark \\
B & Euler 最小密度/最小压强/质量残差 & rho>0,\allowbreak{} p>0,\allowbreak{} |$\Delta$M| <=\allowbreak{} 1e-12 & min rho=\allowbreak{}0.125000,\allowbreak{} min p=\allowbreak{}0.100000,\allowbreak{} defect=\allowbreak{}1.11e-16 \checkmark \\
B & 周期边界下总质量的最大漂移 & <=\allowbreak{} 1e-13 & $4.441\times 10^{-17}$ \checkmark \\
B & sum(flux\_\allowbreak{}net) 与 F\_\allowbreak{}0 - F\_\allowbreak{}N 的差 & 望远镜求和 -> 0(\allowbreak{}舍入级) & 0 \checkmark \\
B & godunov：M(t)-M(0) 与边界通量累计的最大差 & <=\allowbreak{} 1e-14(\allowbreak{}望远镜求和 + 时间插值误差) & $7.772\times 10^{-16}$ \checkmark \\
B & 特征线交点数组的掩码一致性 & 掩码=\allowbreak{}0 的点不得含有效交点 & masked 越界 0 个，有效但非正 0 个 \checkmark \\
\end{longtable}


\section{它会怎样骗你（本章版）}

\begin{enumerate}
  \item \textbf{``质量不变''不是守恒判据。} 出流边界下守恒格式的质量从 $0.5$ 涨到 $0.8$，
        看起来一点都不``守恒''；而账本残差是 $5.8\times10^{-16}$。
        反过来，周期边界下非守恒格式的质量也可能看起来不变，但它的残差是 $O(1)$。
        本项目的做法：帧里同时给质量曲线、累计边界通量与账本残差，
        \texttt{meta.rel\_defect} 与 \texttt{strings.reading\_note} 写明读法。
  \item \textbf{熵条件不是``可选的严格性''。} 跨声速稀疏波上，朴素迎风通量的
        RH 满足、TV 不增、格式单调，所有常规检查都通过；只有熵产从负变正。
        本项目的做法：\texttt{meta.entropy\_production} 进帧，并在图注里把它与
        TV、RH 并列写出；任何``误差很小''的结论必须同时看熵号。
  \item \textbf{破散之后的精确解不是精确解。} 特征回溯给出多值解的一支，
        它光滑、能满足一堆不变量，拿它算误差会得到漂亮的假收敛。
        本项目的做法：\texttt{exact\_valid} 掩码逐点进帧，前端据此把精确解曲线
        截断并把误差轴整条置灰；\texttt{strings.mask\_note} 写明``这是没有参照物，
        不是误差变大''。
  \item \textbf{守恒与精度是两本账。} MUSCL 的残差比一阶大 2 倍却仍然在舍入级；
        二阶重构也不会自动修好边界记账。本项目的做法：账本残差按每个格式单独给，
        而不是拿``整体看起来守恒''一句话代替。
  \item \textbf{CFL 越界不是``误差大一点''。} $\nu>1$ 时格式不再是一个离散解：
        $\max|u|$ 指数增长、通量越界。本项目的做法：帧把解截到 $\pm2.5$ 并置
        \texttt{overflow} 与 \texttt{u\_clipped} 掩码——把失控标出来，而不是藏起来。
\end{enumerate}
